\clearpage
\section{Introdução e preparação do ambiente}
\label{sec:ambiente}

O guia foi preparado como material de apoio para o estudo de Ciência de Dados. Reúne conceitos de Python e Estatística, exemplos de análise de tabelas e gráficos e uma seção de consulta rápida.

\subsection{Como usar os exemplos}

O primeiro passo é ter um ambiente em que o código Python possa ser executado. O Google Colab funciona pelo navegador. Uma instalação local permite executar os exemplos no próprio computador.

Um \emph{notebook} reúne células de código, texto e resultados. Os blocos de Python podem ser executados em ordem no mesmo notebook. Os blocos marcados como PowerShell, Linux ou macOS e Terminal são comandos de preparação para o terminal correspondente.

A leitura de \codigo{notas.csv} na Seção \ref{sec:ler} cria a tabela \codigo{df}. A Seção \ref{sec:ausentes} cria \codigo{df\_validas}, usada nos gráficos com notas. Ao iniciar diretamente pelos gráficos ou reiniciar a sessão, é necessário executar antes as células de leitura e preparação dos dados.

\subsection{Google Colab}

O \href{https://colab.research.google.com/}{Google Colab} permite criar e executar notebooks no navegador. O bloco abaixo pode ser executado em uma célula de código. A saída esperada é o número \codigo{5}, seguido da versão do pandas instalada na sessão.

\begin{exemplo}
import numpy as np
import pandas as pd
import matplotlib.pyplot as plt

print(2 + 3)           # 5
print(pd.__version__)  # Versao do pandas no ambiente
\end{exemplo}

Os exemplos com tabelas e gráficos usam \codigo{notas.csv}. O link \baixarcsv{} fica acessível ao público quando o repositório for aberto. O arquivo também está no painel \textbf{Anexos} do leitor de PDF e na pasta \codigo{dados/} do pacote de distribuição. No Colab, o envio é feito pelo painel \textbf{Arquivos} a cada nova sessão.

\nota{Dependências em notebooks}{Uma biblioteca ausente pode ser instalada em uma célula com \codigo{\%pip install} seguido do nome do pacote. Por exemplo, \codigo{\%pip install scipy} instala a biblioteca usada na seção de probabilidade.}

\subsection{Python e Jupyter no computador}

A execução local começa com a instalação do Python \cite{python}. O ambiente de referência usa Python 3.14.7. O VS Code pode executar scripts \codigo{.py} e notebooks \codigo{.ipynb} com as extensões Python e Jupyter. Os comandos \verb|py --version| no Windows e \verb|python3 --version| no Linux ou macOS mostram a versão instalada.

Na pasta do guia, um ambiente virtual mantém as bibliotecas do projeto separadas das de outras instalações. No Windows, os comandos abaixo criam e ativam esse ambiente no PowerShell.

\begin{exemplo}[language={},title={Windows (PowerShell)}]
py -m venv .venv
.\.venv\Scripts\Activate.ps1
\end{exemplo}

No Prompt de Comando, a ativação é feita com \path{.venv\Scripts\activate.bat}. No Linux ou macOS, a criação e a ativação são feitas com os comandos a seguir.

\begin{exemplo}[language={},title={Linux ou macOS}]
python3 -m venv .venv
source .venv/bin/activate
\end{exemplo}

Com o ambiente ativado, o comando abaixo instala as versões de referência das bibliotecas a partir de \codigo{requirements.txt}, incluído no pacote de distribuição.

\begin{exemplo}[language={},title={Terminal}]
python -m pip install -r requirements.txt
\end{exemplo}

No VS Code, o interpretador deve apontar para a pasta \codigo{.venv}. Em notebooks, o \emph{kernel} é o processo que executa as células e deve usar esse mesmo ambiente. O comando \codigo{jupyter notebook}, executado com o ambiente ativo, abre o Jupyter no navegador \cite{jupyter}.

Ao iniciar o notebook na pasta do guia, o arquivo de dados fica acessível pelo caminho relativo \codigo{dados/notas.csv}. A leitura é feita na Seção \ref{sec:ler}.

\subsection{Erros comuns de preparação}

Ao ocorrer um erro, a última linha da mensagem geralmente indica seu tipo e a informação necessária para iniciar a correção.

\noindent\begin{tabularx}{\linewidth}{@{}lX@{}}
\toprule
\textbf{Mensagem} & \textbf{Causa e como corrigir} \\
\midrule
\codigo{ModuleNotFoundError} & O módulo não foi encontrado. Para instalar a SciPy, por exemplo, execute \codigo{\%pip install scipy} em uma célula do notebook. No computador, ative o ambiente virtual e repita a instalação da Seção 1.3. Consulte o \href{https://packaging.python.org/en/latest/guides/installing-using-pip-and-virtual-environments/}{guia de instalação de pacotes}. \\
\addlinespace
\codigo{FileNotFoundError} & Confira o nome e o caminho do arquivo. No Colab, envie o CSV pelo painel \textbf{Arquivos}. No computador, abra o notebook na pasta do guia e use \codigo{dados/notas.csv}. Para verificar a pasta atual, importe \codigo{Path} com \codigo{from pathlib import Path} e execute \codigo{print(Path.cwd())}. Consulte \href{https://docs.python.org/pt-br/3.14/library/pathlib.html}{caminhos de arquivos}. \\
\addlinespace
\codigo{NameError} & O nome usado ainda não foi definido ou está escrito incorretamente. Confira a grafia (inclusive maiúsculas e minúsculas) e execute primeiro a célula que define a variável. Após reiniciar a sessão, execute as células anteriores em ordem. Consulte \href{https://docs.python.org/pt-br/3.14/tutorial/errors.html}{erros e exceções em Python}. \\
\bottomrule
\end{tabularx}

\nota{Se o erro persistir}{Pesquise o texto completo da mensagem nas documentações indicadas. No terminal, \codigo{python -m pip show scipy} verifica a instalação da SciPy. Se estiver instalada e o erro continuar no notebook, confira o kernel selecionado. Consulte \href{https://jupyter-notebook.readthedocs.io/en/stable/troubleshooting.html}{as soluções de problemas do Jupyter}.}
